\scanpage{67}
% letterhead
\begin{center}
INSTITUT\\
DES HAUTES ÉTUDES SCIENTIFIQUES\\
\rule{1.15cm}{0.4pt}
\end{center}

\uline{Structure Main Theorem}

a) Démonstration par voie globale, (cas $Y$ loc. \uline{noethérien}\deleted{\unsure{)|\S{}}}, $f$ quasi-projectif -- ou du moins compactifiable --

b) Extension au cas $f$ quasi-fini séparé \uline{de prés. finie}, et $Y$ qu. \unsure{cpct|cpt} qu. sép.

On \unsure{prouve} que la question est alors locale sur $\unsure{Y|X}$ (\uline{IV} 8.12.3), cela ramène au cas noethérien. On \unsure{prouve} de plus que l'hypothèse faite implique $f$ quasi-affine (\uline{IV} 8.11.2.) donc on est sous les conditions de a).

Élimination de la méthode globale : par réduction au cas $Y$ \inserted{local ess.} de type fini sur $\operatorname{Spec}\mathbb{Z}$, et passage au complété et récurrence sur la dimension : On montre que le passage au complété est compatible avec l'opération de fermeture intégrale (\uline{IV} 6.14.4.)
% archive sheet number: 66
